\section{Introduction}
Modern smart grid infrastructures demand predictive rather than reactive load management to mitigate the frequency of blackouts and load-shedding events. In developing nations such as Bangladesh, sudden power outages severely disrupt regional distribution networks when ambient temperatures climb or industrial demand surges without warning, overloading feeder transformers \cite{bpdb2023report}. Standard grid management paradigms remain largely reactive, lacking automated early warnings prior to circuit breaker trips.

To address this challenge, we developed \textbf{SmartGrid Sentinel}, an intelligent localized 2-hour ahead load-shedding risk forecasting framework tailored to 143 Upazilas across 15 districts within the Sylhet and Chattogram divisions of Bangladesh. Using a 10-hour historical telemetry lookback (5 sequence steps at 2-hour sample intervals), our system formulates a multi-class sequential prediction task to classify future grid risk into \texttt{Low}, \texttt{Medium}, or \texttt{High} severity levels at forecast horizon $T+2\text{h}$.

We evaluate twelve machine learning and deep temporal architectures, including classical tree ensembles (Random Forest, XGBoost, Decision Trees, Logistic Regression), recurrent networks (LSTM, GRU, CNN-LSTM), and modern sequence Transformers (Autoformer, Informer, PatchTST, Temporal Fusion Transformer, and TimesNet). Furthermore, because deep sequence models are often criticized as black-box systems, SmartGrid Sentinel integrates an intelligent physical rule engine that converts Softmax confidence probabilities and telemetry into human-readable operator advisories and consumer action checklists.

The primary contributions of this work are summarized as follows:
\begin{enumerate}
    \item \textbf{Localized Upazila-Level Risk Forecasting:} Development of a localized 2-hour ahead risk forecasting engine predicting tri-level load-shedding risk across 143 Upazilas and 1,000 distinct feeder streams in Bangladesh.
    \item \textbf{Leakage-Free Comparative Benchmark:} Evaluation of 12 machine learning and temporal deep-learning models under strict chronological train-test splits, preventing temporal data leakage.
    \item \textbf{Comprehensive Empirical Audit:} Inclusion of paired hypothesis testing (McNemar's test and Wilcoxon signed-rank test), controlled feature ablation, and training time complexity analysis to evaluate model performance defensibly.
    \item \textbf{Operational Robustness Evaluation:} Evaluation of model stability under varying operational regimes, specifically diurnal shifts (Peak vs. Off-Peak hours) and seasonal climatic partitions (Summer, Monsoon, Winter).
    \item \textbf{Integrated Intelligent Alert \& Explanation Layer:} Implementation of a physical rule parser bridging deep sequence confidence scores with control room advisories and citizen checklists, moving distribution systems from reactive load shedding to predictive contingency planning.
\end{enumerate}
